% !PRINTED 579
\weberpage{579}

Setzen wir noch ${}^1/_2 (n - 1) = 2l$ oder $= 2l + 1$, je nachdem
$n = 4l + 1$ oder $4l + 3$ ist, so erhalten wir jetzt mit
genauer Vorzeichenbestimmung

\weq{14}{\begin{aligned} A - B &= \left( \frac{k}{n} \right) \sqrt{n}, & n \equiv 1 \ (\mathrm{mod}\ 4) \\ &= i \left( \frac{k}{n} \right) \sqrt{n}, & n \equiv 3 \ (\mathrm{mod}\ 4). \end{aligned}}

\begin{center}
\S.\,172.
\end{center}

\begin{center}
\gesperrt{Die Perioden von} ${}^1/_3 (n - 1)$ \gesperrt{und} ${}^1/_4 (n - 1)$ \gesperrt{Gliedern.}
\end{center}

Wir gehen jetzt zu dem Falle $e = 3$ über, wobei $n - 1$
durch 3 theilbar angenommen werden muss, also $n = 7, 13, 19,
31, 37, 43 \ldots$ Wir bezeichnen mit $\varrho$ eine imaginäre dritte
Einheitswurzel:

\[ \varrho = \frac{-1 + \sqrt{-3}}{2}, \]

und bestimmen die drei Perioden $\eta, \eta_1, \eta_2$ von je ${}^1/_3 (n - 1)$
Gliedern, die, wie wir in \S.\,168 gesehen haben, die Wurzeln
einer cyklischen cubischen Gleichung sind. Nach \S.\,169 ist

\[ \psi_1(\varrho) = \overset{t}{\underset{1,\, n - 2}{\Sigma}} \varrho^{\mathrm{ind}\, t - 2\, \mathrm{ind}(t + 1)}, \]

wofür, da $-2 \equiv 1 \ (\mathrm{mod}\ 3)$, auch

\weq{1}{\psi_1(\varrho) = \overset{t}{\underset{1,\, n - 2}{\Sigma}} \varrho^{\mathrm{ind}(t^2 + t)}}

gesetzt werden kann. Diese Zahl $\psi_1(\varrho)$ kann in einer der beiden
Formen

\weq{2}{\psi_1(\varrho) = a + b\varrho = \frac{A + b\sqrt{-3}}{2}}

dargestellt werden, worin $a, b, A$ ganze Zahlen sind und $A = 2a - b$
ist.

Es ist dann

\weq{3}{\psi_1(\varrho^2) = a + b\varrho^2 = \frac{A - b\sqrt{-3}}{2},}

und die Formeln (10), (15) des \S.\,169 ergeben

\weq{4}{\eta + \eta_1 + \eta_2 = -1}

\weq{5}{(\varrho, \eta)^3 = n \psi_1(\varrho), \quad (\varrho^2, \eta)^3 = n \psi_1(\varrho^2)}
